\section{投资委员会结论}

本轮研究不再只回答“哪里还没启动”，而是将机会拆成两条曲线。第一条是\textbf{静默布局}：资金持续进入、价格仍在低位、盈利或现金流开始改善。第二条是\textbf{已启动高空间}：趋势已确认，但2026--2028收入、利润和EPS扩张仍能为当前价格提供至少约20\%的审慎空间。前者解决安全边际，后者解决机会成本。

\begin{dashboardbox}[A篮子：静默布局]
\begin{tabularx}{\textwidth}{L{2.0cm}L{2.2cm}L{1.5cm}L{1.5cm}L{1.5cm}X}
\toprule
\textbf{排序} & \textbf{标的} & \textbf{现价} & \textbf{目标} & \textbf{空间} & \textbf{动作与定位} \\
\midrule
1 & 渝农商行 601077 & 6.30 & 8.01 & +27.1\% & 核心观察/回撤配置；息差和营收改善最扎实 \\
2 & 徐工机械 000425 & 8.43 & 11.08 & +31.4\% & 回撤配置/业绩验证；赔率最高但周期波动更大 \\
3 & 沪农商行 601825 & 7.66 & 9.69 & +26.5\% & 收益型回撤配置；高股息、低杠杆、盈利增速较弱 \\
\bottomrule
\end{tabularx}
\end{dashboardbox}

\begin{dashboardbox}[B篮子：已启动但盈利空间未走完]
\begin{tabularx}{\textwidth}{L{2.2cm}L{2.1cm}L{1.6cm}L{1.6cm}L{1.6cm}X}
\toprule
\textbf{排序} & \textbf{标的} & \textbf{现价} & \textbf{目标} & \textbf{空间} & \textbf{动作与定位} \\
\midrule
1 & 江波龙 301308 & 587.60 & 802.30 & +36.5\% & 高弹性回撤配置；H1预告验证盈利，严控周期与现金流 \\
2 & 恒瑞医药 600276 & 55.75 & 74.91 & +34.4\% & 趋势回撤核心；创新药、BD与SOTP共同驱动 \\
3 & 工业富联 601138 & 66.27 & 82.29 & +24.2\% & 预告后业绩核心；AI服务器、800G交换机与Q2加速 \\
4 & 中兴通讯 000063 & 40.53 & 48.75 & +20.3\% & 启动后回踩确认；算力占比提升，等待H2利润修复 \\
\bottomrule
\end{tabularx}
\end{dashboardbox}

\section{板块阶段重新分类}

\begin{exhibitbox}[Exhibit 1：当前板块阶段——只有银行仍是一级行业层面的静默吸纳]
\begin{tabularx}{\textwidth}{L{3.0cm}L{2.4cm}L{2.4cm}X}
\toprule
\textbf{阶段} & \textbf{代表方向} & \textbf{资金/价格信号} & \textbf{投资含义} \\
\midrule
静默吸纳 & 银行 & 周净流入32.25亿元；周涨1.89\%；52周位置13\% & 低位、低波动、未形成趋势扩散 \\
启动且可估值 & AI网络、AI存储、创新药 & 核心公司趋势已确认，2026E盈利桥仍给出20\%+空间 & 只做回撤与业绩确认，不追情绪峰值 \\
启动但条件观察 & 国防军工、传媒、先进封装 & 军工和传媒事件启动；封装核心价格已大涨 & 等利润桥、估值或回撤门槛补齐 \\
资金观察 & 计算机 & 周净流入98.31亿元、周涨2.74\%，估值中位 & 有承接但行情已经开始反应 \\
单日反弹 & 医药、汽车、公用事业、机械设备 & 当日净流入、全周仍净流出 & 不可用一天资金替代持续性证据 \\
低位无资金 & 社会服务、美容护理、建筑装饰、非银 & 位置低但周度资金仍流出 & 可能是价值陷阱，不进入核心池 \\
高位资金孤岛 & 半导体设备个股 & 中微/富创连续流入，但一年位置约90\%/82\% & 资金仍在，位置不再低；等估值与价格消化 \\
\bottomrule
\end{tabularx}
\sourcenote{申万研究、乐咕乐股、证券时报数据宝、AStock测算；截止2026-07-10。}
\end{exhibitbox}

\section{为什么是这三只静默核心}

\begin{itemize}
  \item \textbf{渝农商行}：连续8日主力净流入1.46亿元，期间仅涨3.11\%；2026Q1营收同比+8.39\%、归母净利+5.07\%，净息差回升至1.69\%，不良率降至1.07\%。这是资金、基本面和估值最完整的组合。
  \item \textbf{徐工机械}：连续5日净流入2.60亿元，期间仅涨0.12\%；股价位于一年区间约20\%，2026Q1经营现金流同比+153.12\%。市场看到的是汇率拖累后的低利润增速，资金可能在交易海外、矿机和现金流的中期改善。
  \item \textbf{沪农商行}：连续10日净流入1.42亿元，期间反而下跌3.04\%；融资余额仅占流通市值约0.18\%。但2026E利润增速仅约1\%，因此评级低于渝农商行，属于防御型收益资产而非高弹性反转。
\end{itemize}

\section{为什么已启动篮子选择这四只}

\begin{itemize}
  \item \textbf{江波龙}：2026H1官方预告收入220--250亿元、归母净利润92--110亿元，盈利变化已经从主题变成报表事实。按全年200亿元净利、EPS 47.28元重建后，当前价对应约12.4倍PE；空间大，但经营现金流、库存成本层和存储价格反转是硬风险。
  \item \textbf{恒瑞医药}：2026Q1创新药销售45.26亿元、同比+25.75\%，占药品销售61.69\%；华泰预计2026--2028净利94/118/150亿元。当前价格与华泰87.14元、高盛79.20元SOTP锚之间仍有明显差距。
  \item \textbf{工业富联}：H1预告归母净利中值239亿元，隐含Q2约133.05亿元、环比+25.6\%；AI服务器收入同比+230\%以上、800G以上交换机出货同比+140\%。华泰预告后给出EPS 2.82元和93元目标，AStock保守多锚目标82.29元。
  \item \textbf{中兴通讯}：算力产品收入占比已升至27\%，周度资金与放量突破确认启动。当前价约27倍2026E PE，AStock使用比旧卖方更保守的71亿元净利，仍有约20\%空间；但Q1利润和现金流承压，必须等待H2修复。
\end{itemize}

\section{本轮扩展：不再用7只模型代替板块研究}

报告另建54只核心/卫星池、16只H1预告质量表和28只优先研报穿透包。紫光股份、香农芯创、兆易创新、方正科技和中船防务进入预告后重估池；浪潮信息和锐捷网络虽有预告后新研报，但价格已充分反映，列为业绩兑现后的回撤观察；星网宇达因扣非仍亏且投资收益主导扭亏而明确排除。

\begin{riskbox}[最重要的反证]
“连续主力净流入”不等于机构建仓确认，“趋势已经启动”也不等于可以无条件追涨。若银行净息差恶化、徐工现金流不能延续、中兴H2利润不修复、江波龙现金流/库存恶化、恒瑞创新药增长跌破20\%，或工业富联H2净利低于321亿元且毛利率跌破7\%，对应篮子必须降级。
\end{riskbox}
